\documentclass[../main]{subfiles}
\begin{document}

\chapter{Cálculo de Recipientes Sometidos a Presión}

\section{Cálculo de Recipientes Sometidos a Presión}

\subsection{Introducción}

El cálculo de recipientes sometidos a presión es fundamental en la ingeniería y en diversas aplicaciones industriales. Estos recipientes, como calderas, tanques y tuberías, deben ser diseñados para soportar las presiones internas sin fallos estructurales. A continuación, se presentan las fórmulas y conceptos básicos para el cálculo de recipientes sometidos a presión, así como un ejemplo de aplicación.

\info{Recipientes a Presión}{
Un recipiente a presión es un contenedor diseñado para mantener gases o líquidos a una presión significativamente diferente de la presión ambiental. Estos recipientes están sujetos a normas y códigos de diseño rigurosos para garantizar su seguridad y funcionalidad.}

\subsection{Fórmulas Básicas}

En el diseño y análisis de recipientes a presión, se utilizan varias fórmulas esenciales:

1. \textbf{Espesor de pared del recipiente cilíndrico (fórmula de Barlow)}:
    \begin{equation}
    t = \dfrac{P \cdot D}{2 \cdot \sigma}
    \end{equation}
    Donde:
    \begin{itemize}
        \item $t$: Espesor de la pared del recipiente [m].
        \item $P$: Presión interna del recipiente [Pa].
        \item $D$: Diámetro interno del recipiente [m].
        \item $\sigma$: Tensión admisible del material [Pa].
    \end{itemize}

2. \textbf{Fórmula de la presión máxima de trabajo admisible (PMTA)}:
    \begin{equation}
    P = \dfrac{2 \cdot \sigma \cdot t}{D}
    \end{equation}
    Donde:
    \begin{itemize}
        \item $P$: Presión máxima de trabajo admisible [Pa].
        \item $\sigma$: Tensión admisible del material [Pa].
        \item $t$: Espesor de la pared del recipiente [m].
        \item $D$: Diámetro interno del recipiente [m].
    \end{itemize}

3. \textbf{Volumen de un cilindro}:
    \begin{equation}
    V = \pi \cdot r^2 \cdot h
    \end{equation}
    Donde:
    \begin{itemize}
        \item $V$: Volumen del cilindro [m^3].
        \item $r$: Radio del cilindro [m].
        \item $h$: Altura del cilindro [m].
    \end{itemize}
}

\ejemplo{Ejemplo de Cálculo}{
Calcular el espesor de la pared de un recipiente cilíndrico con un diámetro interno de 1.2 m, una presión interna de 5 MPa y una tensión admisible del material de 250 MPa.

Datos:

$P = 5 \cdot 10^6 \, Pa$

$D = 1.2 \, m$

$\sigma = 250 \cdot 10^6 \, Pa$

\begin{equation}
t = \dfrac{P \cdot D}{2 \cdot \sigma} = \dfrac{5 \cdot 10^6 \, Pa \cdot 1.2 \, m}{2 \cdot 250 \cdot 10^6 \, Pa} = 0.012 \, m
\end{equation}
}

\actividad{Responder el siguiente cuestionario}
\begin{enumerate}
\item ¿Qué es un recipiente a presión?
\item ¿Por qué es importante calcular el espesor de la pared de un recipiente sometido a presión?
\item ¿Cuál es la fórmula de Barlow y para qué se utiliza?
\item ¿Qué factores determinan la presión máxima de trabajo admisible (PMTA) de un recipiente?
\item ¿Qué unidades se utilizan para medir la presión interna de un recipiente?
\item ¿Cómo influye el diámetro del recipiente en el espesor necesario de su pared?
\item ¿Qué es la tensión admisible de un material y por qué es relevante en el diseño de recipientes a presión?
\item ¿Cómo se calcula el volumen de un cilindro y por qué es importante en el contexto de recipientes a presión?
\item ¿Qué normas o códigos suelen aplicarse en el diseño de recipientes a presión?
\item ¿Qué podría suceder si un recipiente a presión no se diseña correctamente?
\end{enumerate}

\actividad{Resolver los siguientes problemas}
\begin{enumerate}
\item Calcular el espesor de la pared de un recipiente cilíndrico con un diámetro interno de 0.8 m, una presión interna de 3 MPa y una tensión admisible del material de 200 MPa.
\item Determinar la presión máxima de trabajo admisible (PMTA) para un recipiente con un espesor de pared de 0.015 m, un diámetro interno de 1 m y una tensión admisible del material de 150 MPa.
\item Calcular el volumen de un cilindro con un radio de 0.5 m y una altura de 2 m.
\item Si un recipiente tiene un diámetro interno de 1.5 m y una presión interna de 4 MPa, ¿cuál debe ser el espesor de la pared si la tensión admisible del material es de 300 MPa?
\item Determinar la presión interna que soporta un recipiente con un espesor de pared de 0.01 m, un diámetro interno de 0.5 m y una tensión admisible del material de 250 MPa.
\item Calcular el volumen de un cilindro con un diámetro de 1 m y una altura de 3 m.
\item Si un recipiente tiene un espesor de pared de 0.02 m, un diámetro interno de 0.7 m y soporta una presión interna de 5 MPa, ¿cuál es la tensión admisible del material?
\item Determinar el espesor necesario para un recipiente de 2 m de diámetro interno, soportando una presión de 6 MPa con una tensión admisible del material de 400 MPa.
\item Calcular el volumen de un cilindro con un radio de 0.3 m y una altura de 4 m.
\item Si un recipiente tiene una tensión admisible del material de 200 MPa y un diámetro interno de 1 m, ¿cuál debe ser el espesor de la pared para soportar una presión interna de 2.5 MPa?
\end{enumerate}

\end{document}
